\subsection{2.1. Fundamentos blockchain y contratos
inteligentes}\label{fundamentos-blockchain-y-contratos-inteligentes}

La tecnología~\textbf{\emph{blockchain}}~se define como un sistema de
registro distribuido (\emph{distributed ledger}) que permite almacenar y
gestionar información de forma descentralizada, segura y resistente a
manipulaciones, sin depender de una autoridad central de confianza. Este
paradigma fue introducido por~Satoshi Nakamoto~en 2008, en el contexto
del diseño de Bitcoin, donde se propone un sistema de dinero
electrónico~\emph{peer-to-peer}~basado en un libro mayor compartido
entre múltiples nodos de la red. {[}1{]}

En una blockchain, las transacciones se agrupan en estructuras
denominadas \textbf{bloques}, que se enlazan secuencialmente mediante
\textbf{funciones hash criptográficas}. Cada bloque contiene, entre
otros elementos, un conjunto de transacciones validadas, una marca
temporal y el hash del bloque anterior. Este encadenamiento garantiza la
integridad del sistema, ya que cualquier modificación en un bloque
previamente confirmado alteraría su hash, rompiendo la continuidad de la
cadena y permitiendo detectar manipulaciones de forma inmediata. {[}2{]}

Figura 1. Cadena genérica de bloques

Fuente: Blockchain Technology Overview, NISTIR 8202.

Adicionalmente, blockchain emplea criptografía asimétrica para asegurar
la autenticidad e integridad de las transacciones. Cada participante
dispone de un par de claves criptográficas (clave pública y clave
privada), donde la clave privada se utiliza para firmar digitalmente las
transacciones, y la clave pública permite a otros nodos verificar su
validez. Este mecanismo elimina la necesidad de intermediarios de
confianza, trasladando la verificación al propio sistema distribuido.
{[}2{]}

\subsubsection{2.1.1. ~Tipos de redes
blockchain}\label{tipos-de-redes-blockchain}

Las redes blockchain pueden clasificarse en función de su modelo de
acceso y gobernanza, distinguiéndose principalmente entre redes públicas
(\emph{permissionless}) y redes privadas (\emph{permissioned}).

Las~\textbf{\emph{blockchains}} \emph{\textbf{permisionless}}~permiten
la participación abierta de cualquier usuario, que puede leer el estado
del \emph{ledger} y, dependiendo del protocolo de consenso, participar
en la validación de bloques. Este modelo prioriza la descentralización,
la transparencia y la resistencia a la censura. Ejemplos representativos
incluyen~Bitcoin~y~Ethereum.

Por otro lado, las~\textbf{\emph{blockchains}}
\emph{\textbf{permissioned}}~restringen el acceso a un conjunto
predefinido de participantes autorizados. En este contexto, los nodos
validadores son conocidos y controlados, lo que permite optimizar el
rendimiento, reducir el coste computacional y facilitar el cumplimiento
de requisitos regulatorios. Sin embargo, este modelo introduce un mayor
grado de centralización. Un ejemplo ampliamente adoptado en entornos
empresariales es~Hyperledger Fabric.

En términos generales, las redes públicas priorizan la transparencia y
la descentralización, mientras que las redes privadas favorecen la
eficiencia operativa, el control y la privacidad, lo que explica su
adopción en contextos corporativos.~{[}2{]}

\subsubsection{2.1.2. ~Consenso entre nodos}\label{consenso-entre-nodos}

Uno de los problemas fundamentales en sistemas distribuidos consiste en
lograr que un conjunto de nodos alcance un acuerdo sobre el estado del
sistema, incluso en presencia de fallos o comportamientos maliciosos.
Este problema ha sido ampliamente estudiado en la literatura, destacando
el trabajo de~Leslie Lamport~junto con Shostak y Pease, quienes
formalizaron el consenso en entornos con fallos arbitrarios mediante
el~\emph{Byzantine Generals Problem}, sentando las bases de la
tolerancia a fallos bizantinos (\emph{Byzantine Fault Tolerance}, BFT).
{[}3{]}

En este contexto, un sistema tolerante a fallos bizantinos debe ser
capaz de alcanzar consenso incluso cuando algunos nodos presentan
comportamientos maliciosos o arbitrarios, como el envío de información
falsa o la manipulación deliberada del protocolo. Este modelo resulta
especialmente relevante en redes abiertas y sin permisos, como las
blockchain públicas, donde no existe una relación de confianza previa
entre los participantes.

No obstante, no todos los algoritmos de consenso contemplan este modelo
adversarial. Protocolos clásicos como~\emph{Paxos Protocol}~están
diseñados para entornos con fallos por parada (\emph{crash faults}), en
los que los nodos pueden fallar o dejar de responder, pero no actúan de
forma maliciosa. Esta limitación los hace inadecuados para escenarios
abiertos como blockchain, donde se asume la posible existencia de
actores maliciosos. {[}4{]}

A partir de estos fundamentos, las redes blockchain han desarrollado
mecanismos de consenso específicos que permiten mantener la coherencia
del sistema en entornos descentralizados y potencialmente hostiles.
{[}2{]}Entre los más relevantes destacan:

\begin{itemize}
\item
  \textbf{Proof of Work (PoW)}: basado en la resolución de problemas
  criptográficos que requieren un elevado coste computacional. Este
  mecanismo, utilizado por Bitcoin, proporciona seguridad al hacer
  económicamente inviable la manipulación de la cadena, aunque presenta
  limitaciones en términos de eficiencia energética y escalabilidad.
\item
  \textbf{Proof of Stake (PoS)}: sustituye el esfuerzo computacional por
  la participación económica, donde los validadores son seleccionados en
  función de la cantidad de activos bloqueados (\emph{stake}). Este
  enfoque mejora la eficiencia energética, pero introduce nuevos
  riesgos, como la concentración de poder o ataques derivados del
  problema~\emph{nothing-at-stake}.
\item
  \textbf{Delegated Proof of Stake (DPoS)}: variante de PoS en la que
  los validadores son elegidos mediante mecanismos de votación. Aumenta
  el rendimiento del sistema, aunque reduce el grado de
  descentralización.
\item
  \textbf{Proof of Authority (PoA)}: modelo en el que un conjunto
  limitado de nodos autorizados valida las transacciones. Se utiliza
  principalmente en blockchains privadas, donde la identidad de los
  participantes es conocida.
\end{itemize}

Cada uno de estos mecanismos implica distintos compromisos entre
seguridad, descentralización y rendimiento, lo que en la literatura se
describe frecuentemente como un conjunto de~\emph{trade-offs}~inherentes
al diseño de sistemas blockchain, comúnmente denominado el ``trilema de
blockchain''. {[}5{]}, {[}6{]}

Desde una perspectiva de ciberseguridad, la elección del protocolo de
consenso no solo condiciona la arquitectura del sistema, sino también su
superficie de ataque. En particular, estos compromisos influyen
directamente en la viabilidad de vectores de explotación como ataques
del 51\%, censura de transacciones o manipulación del orden de ejecución
(\emph{front-running}), especialmente en entornos abiertos y sin
permisos.

\subsubsection{2.1.3. ~Fundamentos de contratos
inteligentes}\label{fundamentos-de-contratos-inteligentes}

Los \textbf{contratos inteligentes} (\emph{smart contracts}) constituyen
una extensión funcional de la tecnología blockchain que permite la
ejecución de lógica programable de forma descentralizada. El concepto
fue introducido por~Nick Szabo~en 1994, quien los definió como ``un
protocolo de transacción informatizado que ejecuta los términos de un
contrato'' {[}7{]}

En la actualidad, los contratos inteligentes se materializan como
programas desplegados en una red blockchain que encapsulan tanto lógica
como estado persistente. Su ejecución es llevada a cabo por los nodos de
la red, los cuales deben alcanzar un resultado determinista e idéntico
para garantizar la coherencia global del sistema. Este requisito impone
restricciones relevantes en el diseño del software, especialmente en lo
relativo al uso de fuentes externas de información o funciones no
deterministas.

Cabe destacar que no todas las plataformas blockchain soportan contratos
inteligentes. Redes como~Bitcoin~incorporan un lenguaje de scripting
limitado, diseñado para priorizar la seguridad y la simplicidad,
mientras que otras plataformas como~Ethereum~permiten la ejecución de
código arbitrario, habilitando así el desarrollo de aplicaciones
complejas sobre la blockchain. {[}8{]}

\subsubsection{2.1.4. Ethereum Virtual
Machine}\label{ethereum-virtual-machine}

La adopción generalizada de contratos inteligentes se consolidó con la
aparición de~Ethereum, que introduce un entorno de ejecución específico
denominado~\textbf{\emph{Ethereum Virtual Machine}~(EVM)}. Esta máquina
virtual permite la ejecución de código de propósito general dentro de un
entorno distribuido y determinista.

El funcionamiento interno de Ethereum se describe formalmente en
el~\emph{Yellow Paper}, elaborado por~Gavin Wood, donde se definen:

\begin{itemize}
\tightlist
\item
  La arquitectura y semántica de la EVM
\item
  El modelo de ejecución de contratos y transacciones
\item
  El sistema de gas como mecanismo de control de recursos
\end{itemize}

La EVM es una máquina virtual basada en pila (\emph{stack-based}), donde
cada operación tiene un coste asociado en gas. Este mecanismo permite
limitar el consumo de recursos computacionales y actúa como protección
frente a ataques de denegación de servicio, al imponer un coste
económico a la ejecución de operaciones complejas o potencialmente
abusivas. {[}9{]}

\subsubsection{2.1.5. Propiedades de los contratos
inteligentes}\label{propiedades-de-los-contratos-inteligentes}

Desde un punto de vista técnico, los contratos inteligentes presentan
una serie de propiedades fundamentales{[}2{]}:

\begin{itemize}
\tightlist
\item
  \textbf{Determinismo}: la ejecución debe producir el mismo resultado
  en todos los nodos, garantizando la coherencia del sistema
  distribuido.
\item
  \textbf{Inmutabilidad}: una vez desplegados, los contratos no pueden
  modificarse directamente, lo que dificulta la corrección de errores.
\item
  \textbf{Transparencia}: en redes públicas, el código y el estado del
  contrato son accesibles, favoreciendo la auditabilidad.
\item
  \textbf{Ejecución descentralizada}: no existe un punto único de
  control, eliminando dependencias de entidades centrales.
\end{itemize}

\subsubsection{2.1.6. Limitaciones y
riesgos}\label{limitaciones-y-riesgos}

A pesar de sus ventajas, los contratos inteligentes presentan
importantes desafíos desde el punto de vista de la ciberseguridad
{[}2{]}:

\begin{itemize}
\tightlist
\item
  \textbf{Vulnerabilidades en el código}: errores de implementación
  pueden derivar en fallos críticos explotables.
\item
  \textbf{Problemas de control de acceso}: una gestión incorrecta de
  permisos puede permitir operaciones no autorizadas.
\item
  \textbf{Dependencia de oráculos externos}: introduce confianza en
  terceros y posibles vectores de manipulación.
\item
  \textbf{Inmutabilidad del código}: dificulta la corrección de
  vulnerabilidades tras el despliegue.
\item
  \textbf{Costes de ejecución}: el modelo de gas puede ser explotado
  para provocar condiciones de denegación de servicio.
\end{itemize}

Estas características hacen que los errores en contratos inteligentes
puedan tener consecuencias críticas, como la pérdida irreversible de
activos digitales. {[}2{]}, {[}10{]}

\subsubsection{2.1.7. Aplicaciones descentralizadas
(DApps)}\label{aplicaciones-descentralizadas-dapps}

Sobre la base del modelo de ejecución, las propiedades y las
limitaciones descritas, los contratos inteligentes permiten la
construcción de aplicaciones descentralizadas (\emph{Decentralized
Applications}, DApps), que constituyen sistemas completos en los que la
lógica de negocio se implementa parcial o totalmente mediante contratos
inteligentes. {[}11{]}

A diferencia de las aplicaciones tradicionales basadas en arquitecturas
cliente-servidor, las DApps distribuyen la lógica entre múltiples nodos
de la red, eliminando intermediarios y reduciendo la dependencia de
entidades centralizadas. Este cambio implica una redefinición del modelo
de confianza, donde los usuarios depositan su confianza en el código
ejecutado en la blockchain y en las garantías del protocolo subyacente.

Desde la perspectiva de la ciberseguridad, este paradigma introduce
implicaciones relevantes. La transparencia del código facilita su
auditoría, pero también permite a actores maliciosos analizarlo en busca
de vulnerabilidades. Asimismo, la composición de múltiples contratos y
la dependencia de servicios externos amplían la superficie de ataque,
pudiendo dar lugar a vulnerabilidades complejas en sistemas
interconectados, como ocurre en el ecosistema DeFi.

Además, la inmutabilidad de los contratos implica que los errores no
pueden corregirse fácilmente tras su despliegue, lo que incrementa el
impacto potencial de cualquier vulnerabilidad explotada. {[}12{]}
